\documentclass[]{../../../../tex/classes/styledArticle}
\usepackage{../../../../tex/packages/hebrewSupport}
\usepackage{../../../../tex/packages/mathShortcuts}
\usepackage{../../../../tex/packages/theoremsSupport}


\author{שחר פרץ}
\title{חברה דמוקרטית, תרבות וחינוך $\sim$ \textit{מטלת אמצע קורס}}
\begin{document}
	\maketitle
	
	\textbf{הפרק: }חרדים (עונה 3 פרק 3)
	
	\subsection*{רקע}
	
	במדינת ישראל קיימת האוכלוסיה החרדית המהווה כ־15\% מאוכלוסיית המדינה. מתוכם כ־10\% מגדירים את עצמם כ''חרדים לאומיים``. האוכלוסיה החרדית (הלאומית והלא לאומית) נוטה להתבדל משאר החברה (הדתית והחילונית) הן במקרה החיצוני (לדוגמה, אופן הלבוש), הן במנהגים (לדוגמה, הנוהג לשדרך בגילאים צעירים, בין 16-21) והן במקומות המגורים (מרביתם חיים בקהילות שנמצאים ביישובים חרדים או ברובעים חרדים). 
	
	עם זאת, גם החברה החרדית מחולקת לפלגים בתוך עצמה; ליטאים, ספרדים, ירושלמיים, חסידים וכו'. הזרמים השונים נבדלים מאוד גם בשאר המאפיינים הנזכרים לעיל; בקרב הליטאים מקובל שהנשים מפרנסות את הבית, בעוד בפלגים אחרים לא מקובל שהנשים עובדות בכלל, פלגים שונים מתלבשים מעט שונה, בקרב פלגים מסויימים יש הבדלי מעמדות ברורים, פילוג פוליטי (ספרדים לעומת השאר) והפליה מוסדית, כמו שצוין חלקם חרדים לאומיים שאף משרתים בצה''ל (או לכל הפחות, לא רואים שירות בצה''ל כדבר אסור), וכו'. לכן, לכל כלל והכללה בנוגע לחברה החרדית יש יוצאים מן הכלל. 
	
	ככלל, קהילות חרדיות עוקבות במפורט לאחר ההוראות (הפסיקות) של מנהיגי הקהילה, בדומה לעמדת כוחה של הכנסייה הנוצרית בימים עברו (השוואה ששמעתי בירושלים מחרדי לאומי, שמבקר רבות את הממסד החרדי. דת''לית ליטאית לשעבר שנכחה איתנו בדיון הסכימה איתו). אחת מהנשים בסרטון אמרה על עצמם (החרדים) ש''אנחנו הולכים כמו עדר אחרי מה שאומרים לנו`` (22:12) וש''אני אפילו לא חופרת, לא בודקת, לא שואלת``, זאת בהקשר של חריגה מהנחיות הקורונה כדי ללמוד תורה. מאותו המקרה נוכל ללמוד את חשיבות של ערך לימודי התורה בעבור החברה החרדית, כערך עליון – אחד מהאנשים בסרטון אומר ש''לימוד תורה קודם לחוק`` (22:28). נוסף על כך, חרדים רבים רואים ילודה כפעולה חשובה (''אנחנו מאוד רוצים להביא הרבה ילדים לעולם [...]`` [1:13]). שילוב של שידוכים בגיל מוקדם, רצון ללמוד תורה בערך עליון (על חשבון עבודה, שירות למדינה, וכו'), וריבוי ילדים, מוביל ברוב המקרים לקשיים כלכליים וחוסר רצון להתגייס לצה''ל. סוגיית סירוב הגיוס לצה''ל מעודדת ע''י מרבית הרבנים, שטוענים כי צה''ל לא מאפשר לחרדים ללמוד תורה, וכי הוא עלול לפגוע בחרדים ולגרום לשילובם בחברה. העובדה שמערכת החינוך החרדית מופרדת ממערכת החינוך הדתית לאומית מעודדת את בידול זה. 
	
	כפועל יוצא מכך, חלק ניכר מהחרדים מקבלים יותר (הטבות כלכליות, שירותים סוציאלים וכו') למדינה מאשר מה שהם נותנים (בכסף או בשירות צבאי) למדינה. למרות זאת, מערכות המיסוי והטבות המס בישראל אינם מעודדים עבודה של חרדים (באופן ספציפי), למרות ש''סנקציות כלכליות`` בדמות אימוץ מערכות הטבות מס המצויות במרבית המדינות בעולם (שמתנות קבלת הטבות מסויימות במיצוי יכולת עבודה) הוכחו כיעילות בעידוד עבודה בלי לפגוע (בטווח הארוך) בילודה [מקור לכך אפשר למצוא באחד הפרקים של הפודקאסט של ד''ר עומר מואב. סליחה שאני לא זוכר להפנות לפרק הספציפי]. אי־לכך רבה הביקורת כלפי האוכלוסיה החרדית, שלעיתים מקצינה אף לשנאה של ממש (מספר דוגמאות לכך הובאו בסרטון. לדוגמה, אנשים שצועקים לעבר חרדים דברי נאצה ברחוב). 
	
	עם זאת (הפסקה הזו לא תתקשר לשאר הטקסט, אבל אני חושב שמן הראוי לציינה), יש לציין שבשנים האחרונות החברה החרדית עוברת תהליכי מודרניזציה כלשהם. בסרטון מוצגת הדוגמה בה לכל המרואיינים היה פחות רצון להקים משפחה בהיקף של המשפחה בה הם גדלו (כלומר, משפחות עם 5-7 ילדים במקום 8-12). בפרק של הפודקאסט ''יוצאים בשאלה`` עם יעקב שנקר מובאת דוגמה של אדם שלאחר שחזר בשאלה יצר קשר מחודש עם ההורים והמשפחה שלו לאחר זמן רב, וכאשר הטיח האשמות מסוימות על החברה החרדית ע''פ האופן שהוא גדל בו, אֵחיו סיפרו לו שהמצב איננו כפי שהיה כאשר עזב את החברה החרדית. כך גם ה''חרד''לים``, החרדים הלאומיים, הולכים וגדלים בשיעורם בחברה. 
	
	\subsection*{ניתוח הסוגייה}
	\subsubsection*{מנקודת מבט ליברלית}
	הגישה הליברלית תומכת בין היתר בזכויות אדם ובהבדלים בינהם. לכן היא תיטה לקבל יותר את האוכלוסיה החרדית כאוכלוסיה נפרדת במדינה (אפילו במחיר של היותה מיני־מדינה), ותתנגד להגבלה על חרויות הפרט כגון שירות צבאי (שמגביל את חופש התעסוקה והתנועה של צעירים בגילאים 18-21), כאשר הדבר מתאפשר. עם זאת, תחת ההנחה שהמדינה לא יכולה להתקיים ללא שירות צבאי חובה, הגישה הליברלית תדרוש שוויון באותה הגבלת זכויות, ולכן לא תאפשר שירות של אוכלוסיה ספציפית (חילונית ודתית־לאומית) בעוד שאוכלוסיות אחרות (חרדים, ערבים וכו') לא מחויבות לשירות צבאי. לאור הקבלה של כל צד בחברה כייחודי, הגישה הליברלית לא תתנגד למערכת חינוך נפרדת, ובטח לא לאופן חיים ונורמות חברתיות נפרדות. 
	
	נוסף על כך, הגישה הליברלית מקדמת חשיבה ביקורתית במסגרת ההשתתפות בהליך הדמוקרטי בחברה. לכן, מגוון המגבלות על חופש המידע של האדם החרדי הממוצע יתנגשו על ערכי הגישה הליברלית. כך גם ההליכה ''בעיוורון`` אחרי פסיקות דתיות של קבוצה מצומצמת של אנשים, והפחד מ''מה יחשבו`` על המעשים של כל אחד (שרבים בסרטון גם הביעו: מה יחשבו עלי אם ידעו שהלכתי להתראיין בסליחה על השאלה? איך הפעולות שאני עושה משפיעות על השידוכים שלי? וכולי). 
	
	\subsubsection*{מנקודת מבט רפובליקאית}
	הגישה הרפובליקאית מציבה את אינטרס כלל המדינה מעל האינטרס של הפרט, ורואה חשיבות רבה באופן השתתפותו בחברה. הגישה לפיה התורה מעל החוק סותרת את הגישה הרפובליקאית. באופן דומה אי־שירות בצבא, ותרומה מועטה למדינה במיסים (לאור הכנסות מועטות, ולעיתים אף פשיטת יד; הפלג הירושלמי מקבל הרבה תרומות מחו''ל) לא יתקבלו על ידי אלו בעלי הגישה הרפובליקאית. בעלי הגישה האתנו־לאומית ידרשו אף יותר מתמיכה במנגנוני המדינה – הם ידרשו מהחרדים לראות את עצמם כישראלים ציוניים, וליישר עצמם אידיאולוגית עם התנועה הלאומית הנפוצה במדינה (הציונות). 
	
	בנוסף, בעלי הגישה הרפובליקאית יצפו לבקיאות בנושאי האזרחות המדינה, נושאים שלא נלמדים כלל בחינוך החרדי. 	ככלל, לאור העובדה שמרבית החרדים לא שמים בעדיפות (או מתנגדים באופן כללי) את זהותם הישראלית, הגישה הרפובליקאית תיטה להתנגד להיבטים רבים יותר בחברה החרדית מהגישה הליברלית. 
	
	\subsubsection*{נקודות משותפות}
	נבחין ששתי הגישות יתמכו בגיוס לצה''ל תחת התניות מסוימות, אם כי מסיבות שונות (הגישה הליברלית מטעמי שוויון, בעוד בהגישה הרפובליקאית יותר מטעמים של לראות את האזרח מעורב במהלך המדינה). נקודה משותפת (שגם צוינה במטלה הקודמת) היא ששתי הגישות רואות את החוק כחולש ותקף על כולם. הן הגישה שהתורה קודמת לחוק, והן האוטונומיה של החברה החרדית שבמקרים רבים באה לידי ביטוי בהעדפת בתי דין רבניים על פני מערכת המשפט והדרת המשטרה (מעטים החרדים שיתלוננו למשטרה) מתנגדת לשתי הגישות. 
	
	ב־13:45 מרואיין (שב־6:31 דיבר על כך שהוא לא רואה בעיה בכך שחילוני יחיה כחילוני וחרדי כחרדי) הציג את מדינת ישראל ככשלון משום שהיא איננה כופה את החיים הדתיים ומאפשרת לאדם ''לחיות את אורח החיים הלא מחויב שלו``. מצד אחד אני לא יודע כמה הגישה הזו נפוצה בקרב החרדים, אך מצד שני שמעתי אותה נאמרת מספר פעמים גם מפיהם של דתיים לאומיים. הגישה הליברלית שמדגישה את זכויות הפרט תשלול מדינת הלכה, וכמו כן הגישה האתנו־לאומית שרואה את הלאומיות כערך עליון שמייחד את יושבי המקום, ולא הדת. 
	
	גם החלטת סטנדרטיים שונים כלפי נשים (שמשתנים מאוד בין פלג לפלג) לא קבילים הן בגישה הליברלית שרואה את כל האזרחים של המדינה כשווים בזכויותיהם, והן בגישה הרפובליקאית שלא מבדילה בין אנשים לפי דת, לאום, מין וכיו''ב. 
	
	
	\ndoc
\end{document}